\documentclass[../main-notes.tex]{subfile}

\begin{document}

\justifying

\section{Introduction to the course}



\section{Introduction to Calculus}

%\paragraph{A nonsense that I'm building}
%Sometimes I say that mathematics is a language.
%When we overthink how language can be describe, we can say that have words and those words can be classified as verbs, adjetives, nouns and adverbs.
%And also it have some grammatical rules.
%So, in math, the words and grammatical rules are provided by algebra, meanwhile the thing that help us to classified the mathemtical objectives into groups is calculus.

At some point in the history of humankind there were two open questions: How do you determine the slope of a curve at a single, exact point? And how do you measure the area of a shape with a curved boundary?
The first question was asked by the Greeks.
They were studying lines that only touch a circle once, most commonly known as the tangent line to a circle.
However, for other kinds of graphs it was a mystery, like for a straight line, a parabola, or an ellipse.
After some centuries, the answer to this question led us to the development of the derivative.
On the other hand, the second question led to the creation of the integral.
In a more specific tale, Archimedes was a master of computing areas under parabolas; however, his method was a headache.
Some years later, astronomers like Kepler needed to calculate areas of planetary orbits, and this was solved by the definition of an integral.

In general terms, the derivative allows us to give a quantitative answer to questions about the rate of change at an instant, like ``How fast is it changing right now''?
The integral allows us to give a quantitative answer to questions about the accumulation of change, like ``How much has accumulated in total over time''?
These two questions are fundamentally linked. 
If you have a graph of speed over time, the derivative gives the acceleration at any moment, and the integral gives the total distance traveled.

In this course, we only are going to focus on the first question, ``\textit{How do you determine the slope of a curve at a single, exact point?}''
The answer of this question is going to lead us to the definition of the derivative and that concept it is usefull at GPS navigation, medicine dosing, video games, econocmis, climate models, and other disciplines.

Hence, some of the applications that we are going to explore at the end of the course are related speed, growth, decay, accumulation or in a more abstract concept, flux. 

\paragraph{The big picture}
Calculus is related with arithmetic in the sense that the addition operation finds the total and the substraction operation finds a difference.
In calculus, the integration operation finds a total, meanwhile the derivative finds a rate of change.
The reason of why calculus is more related with arithmetic is because arithmetic give us tools to quantify \textbf{static} amounts, meanwhile calculus give us tools to quantify \textbf{dynamic}, changing amounts.

This is quite different from the math that you learn from previous courses.
In math 3 and math 4 it was more focused on algebra.
And algebra is more like an architect who designs the blueprint, that is the function. 
Calculus is the surveyor and engineer who uses specialized tool to take precises measurements from that blue print.

For example, algebra give us a function that models the height $h(t) = -5t^2 + 20 t +1$.
And can also allow us to answer the question \emph{When does the ball hit the ground?} by solving $h(t) = 0$.
Calculus is going to help us to answer questions like \emph{What is the ball's velocity at the exact instant?} with the derivative and \emph{What is the total distance traveled by the ball between and interval of time?} with integration.

As a wrap up, arithmetic allow us to handle numbers.
Algebra is to handle unknowns and relationships.
And now, calculus will help us to handle change itself.
Finally, the main association that I want you to keep in mind is the following:
Derivatives answers \emph{How fast right now} and the integrals answers \textit{How much in total}.
On the other hand, I also want to make an spoiler.
The integral and derivative are two sides of the same coin or also known as inverse operations.


\section{Recap}

With this in mind, let's recap the main ideas or things from previous math courses that are going to be usefull in this course.


\subsection{Algebra toolkit}

Inverse and composition of functions

\subsection{Type of functions}

\begin{itemize}
    \item Linear
    \item Quadratics 
    \item Cubic
    \item Exponential and Logarithmic functions
    \item Radical functions
    \item Root functions
    \item Two dimensional functions: Conic Sections
\end{itemize}


Types and domain and range.

\subsection{Roots and intersections}

Solve for $x$.
Solve system of equations.

\subsection{Graph of functions}

Use excel.






\paragraph{Ideas to sketch the class}
Show them how to read the math.
To process of demostraions class. 
One to show the procedure to solve the problem.
The other to show what is the graphical representation of the procedure.
The second one is to read the math.


\end{document}
